% Mathematical model for joint sensing- and communication-aware
% underserved-region identification and UAV repositioning.

\subsection{ISAC-assisted underserved-region model}

Consider a set of ground UEs $\mathcal{U}$, terrestrial gNBs
$\mathcal{M}$, and aerial UAV gNBs $\mathcal{Q}$. At control epoch $t$, the
communication subsystem reports the serving-cell RSRP $R_u(t)$, packet-delivery
ratio $P_u(t)$, throughput $C_u(t)$, and delay $D_u(t)$ for UE
$u\in\mathcal{U}$. Using only KPIs recorded by a given experiment, UE $u$ is
classified as underserved when
\begin{equation}
 I_u(t)=\mathbf{1}\!\left\{
 R_u(t)<R_{\min}\ \lor\
 P_u(t)<P_{\min}\ \lor\
 C_u(t)<C_{\min}\ \lor\
 D_u(t)>D_{\max}
 \right\}.
 \label{eq:underserved-ue}
\end{equation}
In the current implementation, $R_u(t)$ is always available; the other terms
are included only when their per-UE traces can be aligned to the same control
epoch.

The ISAC subsystem estimates the UE position rather than determining service
quality. Let $\mathbf{p}_u(t)=[x_u(t),y_u(t)]^T$ be the true UE position used
only as simulation ground truth. The position available to the controller is
\begin{equation}
 \widehat{\mathbf{p}}_u(t)=\mathbf{p}_u(t)+\mathbf{e}_u(t),
 \qquad
 \mathbf{e}_u(t)\sim\mathcal{N}(\mathbf{0},\sigma_s^2\mathbf{I}),
 \label{eq:isac-location}
\end{equation}
where $\sigma_s$ represents ISAC localisation uncertainty. This distinction
prevents perfect ns-3 coordinates from being presented as measured sensing
information.

Partition the service region into grid cells $g\in\mathcal{G}$ with centres
$\mathbf{c}_g$. The sensed UE-density map and underserved-intensity map are
\begin{align}
 \rho_g(t) &= \sum_{u\in\mathcal{U}}
 K_h\!\left(\widehat{\mathbf{p}}_u(t)-\mathbf{c}_g\right), \\
 H_g(t) &= \sum_{u\in\mathcal{U}} I_u(t)
 K_h\!\left(\widehat{\mathbf{p}}_u(t)-\mathbf{c}_g\right),
 \label{eq:underserved-heatmap}
\end{align}
where $K_h(\cdot)$ is a spatial binning or kernel function with bandwidth $h$.
The local underserved share is
\begin{equation}
 Z_g(t)=\frac{H_g(t)}{\rho_g(t)+\epsilon}.
 \label{eq:underserved-share}
\end{equation}
Using both $H_g$ and $Z_g$ distinguishes a dense hotspot containing many weak
UEs from a sparsely populated cell with a high underserved fraction.

For a threshold-based detector, the set of underserved spatial cells is
\begin{equation}
 \widehat{\mathcal{G}}_{\mathrm{US}}(t)=
 \left\{g:H_g(t)\geq H_{\min}\ \land\ Z_g(t)\geq Z_{\min}\right\}.
 \label{eq:threshold-region}
\end{equation}
If a learned detector is evaluated, its feature vector can be written as
\begin{equation}
 \mathbf{x}_g(t)=\left[\rho_g,Z_g,\overline R_g,
 \overline P_g,\overline C_g,\overline D_g,\overline L_g\right]^T,
 \label{eq:grid-features}
\end{equation}
where $\overline L_g$ is serving-cell load. A classifier estimates
$q_g(t)=\Pr(g\in\mathcal{G}_{\mathrm{US}}\mid\mathbf{x}_g(t))$ and declares
cell $g$ underserved when $q_g(t)\geq\tau$. Training and testing must be split
by independent simulation seed or trajectory, not by adjacent grid-time rows.

Given $K\leq|\mathcal{Q}|$ target regions, weighted K-means obtains their
centres by solving
\begin{equation}
 \min_{\boldsymbol{\mu}_1,\ldots,\boldsymbol{\mu}_K}
 \sum_{g\in\widehat{\mathcal{G}}_{\mathrm{US}}}
 H_g(t)\min_{k\in\{1,\ldots,K\}}
 \left\|\mathbf{c}_g-\boldsymbol{\mu}_k\right\|_2^2.
 \label{eq:weighted-clustering}
\end{equation}
The UAV-to-hotspot assignment is
\begin{align}
 \min_{a_{jk}}\quad &
 \sum_{j\in\mathcal{Q}}\sum_{k=1}^{K}
 a_{jk}\left\|\mathbf{q}_j(t)-\boldsymbol{\mu}_k(t)\right\|_2 \\
 \mathrm{s.t.}\quad &
 \sum_k a_{jk}\leq1,\quad
 \sum_j a_{jk}=1,\quad a_{jk}\in\{0,1\},
 \label{eq:uav-assignment}
\end{align}
where $\mathbf{q}_j(t)$ is UAV $j$'s current position. Subject to speed and
area constraints, the selected UAV moves toward its assigned centroid:
\begin{equation}
 \mathbf{q}_j(t+\Delta t)=\mathbf{q}_j(t)+
 \min\!\left(v_{\max}\Delta t,
 \left\|\boldsymbol{\mu}_{k(j)}-\mathbf{q}_j(t)\right\|_2\right)
 \frac{\boldsymbol{\mu}_{k(j)}-\mathbf{q}_j(t)}
 {\left\|\boldsymbol{\mu}_{k(j)}-\mathbf{q}_j(t)\right\|_2}.
 \label{eq:uav-motion}
\end{equation}

Let $\mathcal{G}_{\mathrm{US}}(t)$ be the oracle grid constructed from true UE
positions for evaluation only. Spatial detection precision, recall, and F1 are
computed between $\widehat{\mathcal{G}}_{\mathrm{US}}(t)$ and
$\mathcal{G}_{\mathrm{US}}(t)$. After one-to-one matching between oracle
centroids $\boldsymbol{\mu}_k$ and estimated centroids
$\widehat{\boldsymbol{\mu}}_k$, the hotspot localisation error is
\begin{equation}
 E_{\mathrm{centroid}}=\frac{1}{K}\sum_{k=1}^{K}
 \left\|\boldsymbol{\mu}_k-
 \widehat{\boldsymbol{\mu}}_{\pi(k)}\right\|_2.
 \label{eq:centroid-error}
\end{equation}
The end-to-end value of repositioning is then evaluated using measured network
KPIs, including the fraction of underserved UEs, coverage probability,
fifth-percentile throughput, PDR, and radio-outage time before and after the
control action.
